\chapter{ポートフォリオ構築・リスク・執行の研究}
\label{ch:pcresearch}

\chapref{ch:pc}〜\chapref{ch:cost} で説明した道具が、実際のデータでどれだけ効いたかの研究をまとめる。

\section{推定誤差と単純な方法}

\begin{paperbox}[DeMiguel・Garlappi・Uppal（2009）\cite{dgu2009}]
14 の平均分散型の最適化の方法（ベイズの縮小、制約つき、最小分散など）を、7 つのデータで、すべての資産に等しく投資する $1/N$ のポートフォリオと比べた。標本外のシャープレシオ・確実性等価・回転率のどれでも、一貫して $1/N$ に勝つ方法はなかった。平均分散の最適化が $1/N$ に勝つには、25 資産なら約 3,000 か月、50 資産なら約 6,000 か月の推定の期間が要る、と計算した。
\end{paperbox}

この結果は、期待リターンの推定の誤差がいかに大きいかを示す。ただし、$1/N$ は「予測を持たない」場合の比較であり、予測の情報を持つ運用者には、Black--Litterman のように予測を控えめに混ぜる方法が向く。

\begin{itemize}
\item \textbf{Jagannathan・Ma（2003）}\cite{jagannathanma2003}：最小分散のポートフォリオに空売りの禁止の制約をかけると、標本の共分散を縮めたのと同じ効果が生じ、標本外の成績が良くなる。「間違った制約が役に立つ」。
\item \textbf{López de Prado（2016）}\cite{hrp2016}：階層的リスクパリティ（\chapref{ch:pc}）が、モンテカルロの実験で、二次計画による最小分散より標本外の分散が小さいことを示した。ただし、実際のデータでの優位は、比較の条件に依存するという報告もあり、確立した結果とは言いにくい。
\end{itemize}

\section{共分散の推定}

\begin{itemize}
\item \textbf{Ledoit・Wolf（2004）}\cite{ledoitwolf2004}：標本の共分散を、すべての相関を等しいとした行列に向けて縮めると、標本外の最小分散のポートフォリオの変動率が下がる。縮める強さは、データから解析的に計算できる。
\item \textbf{Ledoit・Wolf（2020）}\cite{ledoitwolf2020}：固有値ごとに違う強さで縮める非線形の縮小を、解析的な式で速く計算する方法を示した。資産の数が多いときに、線形の縮小より良い。
\item \textbf{ランダム行列}：Laloux ほか（1999）\cite{laloux1999} と Plerou ほか（1999）\cite{plerou1999} は、株の相関行列の固有値の大部分が、ランダムな行列の固有値の分布（Marchenko--Pastur）の範囲に収まり、範囲の外の少数の固有値（市場全体、業種）だけが安定した構造を持つことを示した。
\end{itemize}

\section{ボラティリティの予測と、量の調整}

Hansen・Huang・Shek（2012）\cite{hhs2012} の Realized GARCH は、日中の高頻度のデータで測った実現の変動率を GARCH に組み込み、日次のリターンだけの GARCH より良い変動率の予測をした。Corsi（2009）\cite{corsi2009} の HAR モデルとともに、実現の変動率による予測は、日次のデータだけのモデルより一般に良い。

\begin{paperbox}[Moreira・Muir（2017）と、その批判]
Moreira・Muir（2017）\cite{moreiramuir2017} は、市場やファクターのポートフォリオの量を、前月の実現の分散の逆数に比例して調整する（変動が大きいときに減らす）と、大きなアルファが生じ、シャープレシオが上がることを示した。変動率は予測できるが、平均リターンは変動率に比例して上がらないためである。

Cederburg・O'Doherty・Wang・Yan（2020）\cite{cederburg2020} は、100 を超える株式の戦略で同じ調整を検証し、その場その場で実行できる形（調整の大きさを事前に決める）では、調整した戦略が調整しない戦略に一貫して勝つわけではないことを示した。標本の全期間を使って調整の大きさを決めると、成績が良く見える。
\end{paperbox}

モメンタムのような急落する戦略では、量の調整の効果が大きい（Barroso・Santa-Clara 2015 \cite{barroso2015}、第 \ref{ch:equityfactor} 章）。一方、市場全体の量の調整の効果は、検証の仕方に大きく依存する。

\section{取引コストと容量}

\subsection{インパクトのモデル}

\begin{itemize}
\item \textbf{Almgren・Thum・Hauptmann・Li（2005）}\cite{athl2005}：大手の証券会社の米国株の大口の注文のデータから、一時的なインパクトが売買の速さの約 0.6 乗に比例する（平方根に近い）ことなどを推定した。
\item \textbf{Tóth ほか（2011）}\cite{toth2011}：大口の注文のインパクトが、注文の大きさの平方根に比例することを、多くの市場で確かめ、板に表れない潜在的な流動性が、注文の大きさに比例して薄くなるというモデルで説明した。
\item 平方根則の多くの実証は、Bouchaud ほかの教科書 \cite{bouchaud2018} にまとめられている。
\end{itemize}

\subsection{コストの後の収益と容量}

\begin{itemize}
\item \textbf{Korajczyk・Sadka（2004）}\cite{korajczyksadka2004}：インパクトを含む取引コストを入れると、モメンタムの戦略の収益は、運用の資金が大きくなるにつれて小さくなり、ある資金の大きさ（損益分岐の規模）でゼロになることを示した。取引の量を流動性に合わせる工夫をした一部のモメンタムの戦略では、見かけの収益の機会が消えるまでに最大で約 50 億ドル（1999 年 12 月の時価総額の水準で）を投じられるが、よく研究されてきた従来のモメンタムの戦略の多くは、大きな規模では実行できないと推計した。
\item \textbf{Novy-Marx・Velikov（2016）}\cite{novymarxvelikov2016}：回転率の低いアノマリーはコストの後も残るが、高いものは多くが消える（第 \ref{ch:equityfactor} 章）。
\item \textbf{Frazzini・Israel・Moskowitz（2018）}\cite{fim2018}：大手の運用会社の実際の大量の売買の記録から取引コストを推定し、学術研究で使われてきたコストの推定より、実際のコストはかなり小さいこと、そのため割安・モメンタムなどの戦略の容量は、これまでの推計より大きいことを論じた。運用会社自身の研究であり、執行の上手な大手だから実現できるコストである点に注意が要る。
\item \textbf{DeMiguel・Martín-Utrera・Nogales・Uppal（2020）}\cite{dmnr2020}：多くの特性を同時に使うと、それぞれの売買が一部打ち消し合い（ある特性で買い、別の特性で売る）、コストが下がる。そのため、取引コストを考えると、むしろ多くの特性を使うことが有利になる。
\end{itemize}

\section{強化学習と深層学習}

\subsection{執行}

最適執行（\chapref{ch:cost}）は、状態（残りの量、時間、板の状態）に応じて行動（どれだけ売買するか）を決める逐次の意思決定の問題で、強化学習に向く。

\begin{itemize}
\item \textbf{Nevmyvaka・Feng・Kearns（2006）}\cite{nfk2006}：米国株の板のデータを使い、強化学習で指値の位置と量を決める執行の方策を学習し、単純な方策よりコストを下げた。
\item \textbf{Ning・Lin・Jaimungal（2021）}\cite{nlj2021}：深層の Q 学習（二重 DQN）で最適執行を学習し、TWAP などより良い成績を示した。
\item \textbf{Hafsi・Vittori}\cite{hafsivittori}：多数の売買の主体を模した市場のシミュレータ（ABIDES）の中で、板の状態を入力にした DQN の執行の方策を学習し、TWAP などの単純な方策より執行のコストを下げた（Deep Research の答え 2 が挙げた研究で、内容はその要約による。題名・著者は照合で確かめ、2026 年の国際会議 ACDSA の会議録に載っている。シミュレータの中の結果であり、実際の市場での検証ではない）。
\end{itemize}

執行は、データが豊富で（毎日、多くの注文）、結果がすぐに分かり、目的（コスト）がはっきりしているので、強化学習が実用に近い分野とされる。

\subsection{ヘッジ}

Buehler・Gonon・Teichmann・Wood（2019）\cite{deephedging2019} の「深層ヘッジ」は、取引コストなどの摩擦がある市場で、デリバティブのヘッジの方策をニューラルネットで直接学習する枠組みを示した。実際の銀行のデリバティブの業務にも使われている、と報じられている。

\subsection{ポートフォリオの配分}

Jiang・Xu・Liang（2017）\cite{jxl2017} や、オープンソースの FinRL \cite{finrl2020} など、深層の強化学習でポートフォリオの比率を直接決める研究は多い。しかし、

\begin{itemize}
\item 報酬（リターン）の信号が極めて弱く、学習が不安定で過学習しやすい。
\item 市場が変わるので、過去の環境で学んだ方策が将来も良いとは限らない。
\item 取引コストや市場への影響を、シミュレーションで正しく再現しにくい。
\end{itemize}
という理由で、標本外の頑健な成績を示した研究は少なく、実務での利用も限られる（Deep Research の答え 2 も、研究の段階と評価していた）。

\begin{keybox}[まとめ]
\begin{description}[leftmargin=1em,style=nextline]
\item[確立している]
  期待リターンの推定の誤差は非常に大きく、単純な方法（$1/N$、縮小、制約）が強い。共分散の縮小推定と実現の変動率による予測は有効。インパクトは平方根則に近い。コストは戦略の容量を決め、回転率の高い戦略の収益の多くを消す。強化学習は執行では実用に近い。
\item[議論が続いている]
  変動率による量の調整の効果（検証の仕方に依存）。階層的リスクパリティの実データでの優位。実際の取引コストの大きさ（学術の推定と運用会社の推定の差）。強化学習による配分。
\end{description}
\end{keybox}
